\begin{afterword}
\summaryhead

The post is a passage from Robert Pirsig's \textsc{Zen and the Art of Motorcycle Maintenance}, introduced by two sentences: the author does not know whether the story is based on reality, but ``either way, it's true.''

In the passage a writing teacher has students who cannot think of anything to say. One wants to write about the United States; he tells her to narrow it to Bozeman, then to its main street, and still she writes nothing. At last, angry, he tells her to write about the front of one building, the Opera House, starting with the upper left-hand brick. She returns with five thousand words. Thinking it over, he concludes that she had been trying to repeat what she had already heard, and that narrowing the subject to one brick forced her to look and see for herself.

\responsehead

The post is a quotation with a two-sentence frame. The only claim it makes in its own voice is in the frame: the story is true whether or not it happened.

That claim is the weak point. If the story may be invented, it cannot also be the evidence that its lesson is true, and the post offers no other evidence. Two days later, in ``The Logical Fallacy of Generalization from Fictional Evidence,'' the author asked whether an issue does not ``have to be argued in its own right, regardless of how Vinge depicted his characters.'' Here the post argues by story alone. The doubt in the frame is also partly answered by the book itself: Pirsig's Author's Note says the book is ``based on actual occurrences'' and ``must be regarded in its essence as fact,'' although much was changed.

The book gives more support than the post quotes. In the next paragraphs the teacher has whole classes write for an hour about his thumb or a coin, and reports that no one complained of having nothing to say. That is still one teacher's account, but it is the passage's only sign that the method worked more than once.

What the post adds is its place in the book. It follows ``Cached Thoughts,'' and the teacher's diagnosis, that the student was trying to repeat what she had heard, is close to that post's idea. The post does not draw the connection; the reader has to.

In short: a quoted anecdote offered as true ``either way,'' with no evidence beyond the story, although the book itself claims to be fact in essence and gives further examples the post does not quote.
\end{afterword}
